
\documentclass{article}
\usepackage{colortbl}
\usepackage{makecell}
\usepackage{multirow}
\usepackage{supertabular}

\begin{document}

\newcounter{utterance}

\centering \large Interaction Transcript for game `hot\_air\_balloon', experiment `air\_balloon\_survival\_en\_reasoning off\_hard', episode 0 with deepseek{-}chat{-}v3.1{-}t1.0.
\vspace{24pt}

{ \footnotesize  \setcounter{utterance}{1}
\setlength{\tabcolsep}{0pt}
\begin{supertabular}{c@{$\;$}|p{.15\linewidth}@{}p{.15\linewidth}p{.15\linewidth}p{.15\linewidth}p{.15\linewidth}p{.15\linewidth}}
    \# & $\;$A & \multicolumn{4}{c}{Game Master} & $\;\:$B\\
    \hline

    \theutterance \stepcounter{utterance}  
    & & \multicolumn{4}{p{0.6\linewidth}}{
        \cellcolor[rgb]{0.9,0.9,0.9}{
            \makecell[{{p{\linewidth}}}]{
                \texttt{\tiny{[P1$\langle$GM]}}
                \texttt{Stai partecipando a un gioco di negoziazione collaborativa.} \\
\\ 
\texttt{Insieme a un altro partecipante, dovete accordarvi su un unico set di oggetti da conservare. Ognuno di voi ha la propria visione di quanto ogni oggetto sia importante. Non sai come l'altro partecipante valuta gli oggetti. Inoltre, ti viene fornito lo sforzo che ogni oggetto richiede.} \\
\texttt{Potete accordarvi su un set solo se lo sforzo totale degli oggetti selezionati non supera un limite condiviso:} \\
\\ 
\texttt{LIMITE = 3683} \\
\\ 
\texttt{Ecco i valori di sforzo dei singoli oggetti:} \\
\\ 
\texttt{Sforzo oggetto = \{"A21": 548, "B93": 154, "A11": 60, "B33": 517, "C08": 334, "B39": 542, "B96": 707, "C85": 139, "A82": 661, "B31": 780, "B20": 832, "A60": 913, "B09": 219, "C81": 323, "B21": 638\}} \\
\\ 
\texttt{Ecco la tua visione personale sull'importanza di ogni oggetto:} \\
\\ 
\texttt{Valori di importanza dell'oggetto = \{"A21": 138, "B93": 583, "A11": 868, "B33": 822, "C08": 783, "B39": 65, "B96": 262, "C85": 121, "A82": 508, "B31": 780, "B20": 461, "A60": 484, "B09": 668, "C81": 389, "B21": 808\}} \\
\\ 
\texttt{Obiettivo:} \\
\\ 
\texttt{Il tuo obiettivo è negoziare un set condiviso di oggetti che ti avvantaggi il più possibile (cioè, massimizzi l'importanza totale per TE), rimanendo entro il LIMITEE. Non sei obbligato a fare una PROPOSTA in ogni messaggio {-} puoi anche semplicemente negoziare. Tutte le tattiche sono permesse!} \\
\\ 
\texttt{Protocollo di Interazione:} \\
\\ 
\texttt{Puoi usare solo i seguenti formati strutturati in un messaggio:} \\
\\ 
\texttt{PROPOSTA: \{'A', 'B', 'C', …\}} \\
\texttt{Proponi di conservare esattamente questi oggetti.} \\
\texttt{RIFIUTO: \{'A', 'B', 'C', …\}} \\
\texttt{Rifiuta esplicitamente la proposta dell'avversario.} \\
\texttt{ARGOMENTAZIONE: \{'...'\}} \\
\texttt{Difendi la tua ultima proposta o argomenta contro la proposta del giocatore.} \\
\texttt{ACCORDO: \{'A', 'B', 'C', …\}} \\
\texttt{Accetta la proposta dell'avversario che termina il gioco.} \\
\\ 
\\ 
\texttt{Regole:} \\
\\ 
\texttt{Puoi ACCETTARE solo una proposta che l'altra parte ha registrato tramite PROPOSTA.} \\
\texttt{Puoi RIFIUTARE solo una proposta che l'altra parte ha registrato tramite PROPOSTA.} \\
\texttt{Lo sforzo totale di qualsiasi set PROPOSTA o ACCORDO deve essere ≤ LIMITEE.} \\
\texttt{NON rivelare i tuoi punteggi di importanza nascosti.} \\
\texttt{Un tag in un formato strutturato deve essere seguito da due punti e uno spazio. L'argomento deve essere un set python contenente 0 o più stringhe.} \\
\texttt{Quindi, deve essere della forma TAG: \{...\}} \\
\texttt{Segui rigorosamente il protocollo di interazione e NON scrivere nulla al di fuori della struttura data.} \\
\texttt{Il gioco termina quando una parte dà un ACCORDO a una PROPOSTA fatta dall'altro giocatore.} \\
\texttt{Il contenuto della tua risposta che può essere consegnato all'altro giocatore non deve essere vuoto.} \\
\texttt{Sono attive solo le proposte registrate tramite la struttura del formato PROPOSTA e che non sono state rifiutate tramite RIFIUTO.} \\
\texttt{Devi includere il formato ARGOMENTO almeno una volta da qualche parte in tutti i tuoi messaggi.} \\
            }
        }
    }
    & \\ \\

    \theutterance \stepcounter{utterance}  
    & \multicolumn{4}{p{0.6\linewidth}}{
        \cellcolor[rgb]{0.9,0.9,0.9}{
            \makecell[{{p{\linewidth}}}]{
                \texttt{\tiny{[P1$\rangle$GM]}}
                \texttt{PROPOSTA: \{'A11', 'B93', 'B33', 'C08', 'B09', 'C81', 'B21', 'A82', 'B20'\}} \\
\texttt{ARGOMENTAZIONE: \{'Questa proposta include oggetti ad alto valore con sforzi moderati, offrendo un buon equilibrio. Abbiamo A11, B33, C08 e B21 che sono particolarmente preziosi, insieme ad altri che completano bene il set senza superare il limite.'\}} \\
            }
        }
    }
    & & \\ \\

    \theutterance \stepcounter{utterance}  
    & & \multicolumn{4}{p{0.6\linewidth}}{
        \cellcolor[rgb]{0.9,0.9,0.9}{
            \makecell[{{p{\linewidth}}}]{
                \texttt{\tiny{[GM$\rangle$P2]}}
                \texttt{Stai partecipando a un gioco di negoziazione collaborativa.} \\
\\ 
\texttt{Insieme a un altro partecipante, dovete accordarvi su un unico set di oggetti da conservare. Ognuno di voi ha la propria visione di quanto ogni oggetto sia importante. Non sai come l'altro partecipante valuta gli oggetti. Inoltre, ti viene fornito lo sforzo che ogni oggetto richiede.} \\
\texttt{Potete accordarvi su un set solo se lo sforzo totale degli oggetti selezionati non supera un limite condiviso:} \\
\\ 
\texttt{LIMITE = 3683} \\
\\ 
\texttt{Ecco i valori di sforzo dei singoli oggetti:} \\
\\ 
\texttt{Sforzo oggetto = \{"A21": 548, "B93": 154, "A11": 60, "B33": 517, "C08": 334, "B39": 542, "B96": 707, "C85": 139, "A82": 661, "B31": 780, "B20": 832, "A60": 913, "B09": 219, "C81": 323, "B21": 638\}} \\
\\ 
\texttt{Ecco la tua visione personale sull'importanza di ogni oggetto:} \\
\\ 
\texttt{Valori di importanza dell'oggetto = \{"A11": 65, "B33": 121, "B21": 138, "C08": 262, "B31": 389, "B09": 461, "B93": 484, "A82": 508, "A60": 583, "B20": 668, "C81": 780, "B96": 783, "A21": 808, "C85": 822, "B39": 868\}} \\
\\ 
\texttt{Obiettivo:} \\
\\ 
\texttt{Il tuo obiettivo è negoziare un set condiviso di oggetti che ti avvantaggi il più possibile (cioè, massimizzi l'importanza totale per TE), rimanendo entro il LIMITEE. Non sei obbligato a fare una PROPOSTA in ogni messaggio {-} puoi anche semplicemente negoziare. Tutte le tattiche sono permesse!} \\
\\ 
\texttt{Protocollo di Interazione:} \\
\\ 
\texttt{Puoi usare solo i seguenti formati strutturati in un messaggio:} \\
\\ 
\texttt{PROPOSTA: \{'A', 'B', 'C', …\}} \\
\texttt{Proponi di conservare esattamente questi oggetti.} \\
\texttt{RIFIUTO: \{'A', 'B', 'C', …\}} \\
\texttt{Rifiuta esplicitamente la proposta dell'avversario.} \\
\texttt{ARGOMENTAZIONE: \{'...'\}} \\
\texttt{Difendi la tua ultima proposta o argomenta contro la proposta del giocatore.} \\
\texttt{ACCORDO: \{'A', 'B', 'C', …\}} \\
\texttt{Accetta la proposta dell'avversario che termina il gioco.} \\
\\ 
\\ 
\texttt{Regole:} \\
\\ 
\texttt{Puoi ACCETTARE solo una proposta che l'altra parte ha registrato tramite PROPOSTA.} \\
\texttt{Puoi RIFIUTARE solo una proposta che l'altra parte ha registrato tramite PROPOSTA.} \\
\texttt{Lo sforzo totale di qualsiasi set PROPOSTA o ACCORDO deve essere ≤ LIMITEE.} \\
\texttt{NON rivelare i tuoi punteggi di importanza nascosti.} \\
\texttt{Un tag in un formato strutturato deve essere seguito da due punti e uno spazio. L'argomento deve essere un set python contenente 0 o più stringhe.} \\
\texttt{Quindi, deve essere della forma TAG: \{...\}} \\
\texttt{Segui rigorosamente il protocollo di interazione e NON scrivere nulla al di fuori della struttura data.} \\
\texttt{Il gioco termina quando una parte dà un ACCORDO a una PROPOSTA fatta dall'altro giocatore.} \\
\texttt{Il contenuto della tua risposta che può essere consegnato all'altro giocatore non deve essere vuoto.} \\
\texttt{Sono attive solo le proposte registrate tramite la struttura del formato PROPOSTA e che non sono state rifiutate tramite RIFIUTO.} \\
\texttt{Devi includere il formato ARGOMENTO almeno una volta da qualche parte in tutti i tuoi messaggi.} \\
\\ 
\\ 
\texttt{Ora riceverai il primo messaggio dell'altro giocatore.} \\
\\ 
\texttt{PROPOSTA: \{'A11', 'B93', 'B33', 'C08', 'B09', 'C81', 'B21', 'A82', 'B20'\}} \\
\texttt{ARGOMENTAZIONE: \{'Questa proposta include oggetti ad alto valore con sforzi moderati, offrendo un buon equilibrio. Abbiamo A11, B33, C08 e B21 che sono particolarmente preziosi, insieme ad altri che completano bene il set senza superare il limite.'\}} \\
            }
        }
    }
    & \\ \\

    \theutterance \stepcounter{utterance}  
    & & & \multicolumn{4}{p{0.6\linewidth}}{
        \cellcolor[rgb]{0.9,0.9,0.9}{
            \makecell[{{p{\linewidth}}}]{
                \texttt{\tiny{[GM$\langle$P2]}}
                \texttt{PROPOSTA: \{'A11', 'B93', 'B33', 'C08', 'B09', 'C81', 'B21', 'A82', 'A21', 'C85'\}} \\
\texttt{ARGOMENTAZIONE: \{'La tua proposta è interessante ma ho aggiunto A21 e C85 che sono oggetti molto importanti per me, rimuovendo B20 che richiede molto sforzo. Questo set mantiene un buon equilibrio con uno sforzo totale di 3680, appena sotto il limite, e include molti oggetti di valore per entrambi.'\}} \\
            }
        }
    }
     \\ \\

    \theutterance \stepcounter{utterance}  
    & & \multicolumn{4}{p{0.6\linewidth}}{
        \cellcolor[rgb]{0.9,0.9,0.9}{
            \makecell[{{p{\linewidth}}}]{
                \texttt{\tiny{[P1$\langle$GM]}}
                \texttt{PROPOSTA: \{'A11', 'B93', 'B33', 'C08', 'B09', 'C81', 'B21', 'A82', 'A21', 'C85'\}} \\
\texttt{ARGOMENTAZIONE: \{'La tua proposta è interessante ma ho aggiunto A21 e C85 che sono oggetti molto importanti per me, rimuovendo B20 che richiede molto sforzo. Questo set mantiene un buon equilibrio con uno sforzo totale di 3680, appena sotto il limite, e include molti oggetti di valore per entrambi.'\}} \\
            }
        }
    }
    & \\ \\

    \theutterance \stepcounter{utterance}  
    & \multicolumn{4}{p{0.6\linewidth}}{
        \cellcolor[rgb]{0.9,0.9,0.9}{
            \makecell[{{p{\linewidth}}}]{
                \texttt{\tiny{[P1$\rangle$GM]}}
                \texttt{PROPOSTA: \{'A11', 'B93', 'B33', 'C08', 'B09', 'C81', 'B21', 'A82', 'A21'\}} \\
\texttt{ARGOMENTAZIONE: \{'Apprezzo l'aggiunta di A21 nella tua proposta. Tuttavia, C85 ha un valore limitato rispetto al suo sforzo. Rimuovendo C85, possiamo mantenere un set forte con A11, B33, C08, B21 e altri oggetti ad alto valore. Questo set ha uno sforzo totale di 3541, ben sotto il limite, lasciando spazio per potenziali aggiunte future se necessario.'\}} \\
            }
        }
    }
    & & \\ \\

    \theutterance \stepcounter{utterance}  
    & & \multicolumn{4}{p{0.6\linewidth}}{
        \cellcolor[rgb]{0.9,0.9,0.9}{
            \makecell[{{p{\linewidth}}}]{
                \texttt{\tiny{[GM$\rangle$P2]}}
                \texttt{PROPOSTA: \{'A11', 'B93', 'B33', 'C08', 'B09', 'C81', 'B21', 'A82', 'A21'\}} \\
\texttt{ARGOMENTAZIONE: \{'Apprezzo l'aggiunta di A21 nella tua proposta. Tuttavia, C85 ha un valore limitato rispetto al suo sforzo. Rimuovendo C85, possiamo mantenere un set forte con A11, B33, C08, B21 e altri oggetti ad alto valore. Questo set ha uno sforzo totale di 3541, ben sotto il limite, lasciando spazio per potenziali aggiunte future se necessario.'\}} \\
            }
        }
    }
    & \\ \\

    \theutterance \stepcounter{utterance}  
    & & & \multicolumn{4}{p{0.6\linewidth}}{
        \cellcolor[rgb]{0.9,0.9,0.9}{
            \makecell[{{p{\linewidth}}}]{
                \texttt{\tiny{[GM$\langle$P2]}}
                \texttt{PROPOSTA: \{'A11', 'B93', 'B33', 'C08', 'B09', 'C81', 'B21', 'A82', 'A21', 'B31'\}} \\
\texttt{ARGOMENTAZIONE: \{'Accetto di rimuovere C85 ma propongo di aggiungere B31 che offre un eccellente rapporto valore/sforzo. Questo set ha uno sforzo totale di 3679, appena sotto il limite, e include oggetti chiave come A21, B31 e B21 che sono molto preziosi per me. B31 completa bene il set senza compromettere il limite.'\}} \\
            }
        }
    }
     \\ \\

    \theutterance \stepcounter{utterance}  
    & & \multicolumn{4}{p{0.6\linewidth}}{
        \cellcolor[rgb]{0.9,0.9,0.9}{
            \makecell[{{p{\linewidth}}}]{
                \texttt{\tiny{[P1$\langle$GM]}}
                \texttt{PROPOSTA: \{'A11', 'B93', 'B33', 'C08', 'B09', 'C81', 'B21', 'A82', 'A21', 'B31'\}} \\
\texttt{ARGOMENTAZIONE: \{'Accetto di rimuovere C85 ma propongo di aggiungere B31 che offre un eccellente rapporto valore/sforzo. Questo set ha uno sforzo totale di 3679, appena sotto il limite, e include oggetti chiave come A21, B31 e B21 che sono molto preziosi per me. B31 completa bene il set senza compromettere il limite.'\}} \\
            }
        }
    }
    & \\ \\

    \theutterance \stepcounter{utterance}  
    & \multicolumn{4}{p{0.6\linewidth}}{
        \cellcolor[rgb]{0.9,0.9,0.9}{
            \makecell[{{p{\linewidth}}}]{
                \texttt{\tiny{[P1$\rangle$GM]}}
                \texttt{RIFIUTO: \{'A11', 'B93', 'B33', 'C08', 'B09', 'C81', 'B21', 'A82', 'A21', 'B31'\}} \\
\texttt{ARGOMENTAZIONE: \{'B31 richiede uno sforzo molto elevato (780) per il suo valore relativo. Preferirei sostituirlo con oggetti più efficienti come C85 o includere B20 che, nonostante lo sforzo, offre un buon contributo. Il set proposto è troppo vicino al limite senza ottimizzare sufficientemente il valore totale.'\}} \\
            }
        }
    }
    & & \\ \\

    \theutterance \stepcounter{utterance}  
    & & \multicolumn{4}{p{0.6\linewidth}}{
        \cellcolor[rgb]{0.9,0.9,0.9}{
            \makecell[{{p{\linewidth}}}]{
                \texttt{\tiny{[GM$\rangle$P2]}}
                \texttt{RIFIUTO: \{'A11', 'B93', 'B33', 'C08', 'B09', 'C81', 'B21', 'A82', 'A21', 'B31'\}} \\
\texttt{ARGOMENTAZIONE: \{'B31 richiede uno sforzo molto elevato (780) per il suo valore relativo. Preferirei sostituirlo con oggetti più efficienti come C85 o includere B20 che, nonostante lo sforzo, offre un buon contributo. Il set proposto è troppo vicino al limite senza ottimizzare sufficientemente il valore totale.'\}} \\
            }
        }
    }
    & \\ \\

    \theutterance \stepcounter{utterance}  
    & & & \multicolumn{4}{p{0.6\linewidth}}{
        \cellcolor[rgb]{0.9,0.9,0.9}{
            \makecell[{{p{\linewidth}}}]{
                \texttt{\tiny{[GM$\langle$P2]}}
                \texttt{PROPOSTA: \{'A11', 'B93', 'B33', 'C08', 'B09', 'C81', 'B21', 'A82', 'A21', 'B39'\}} \\
\texttt{ARGOMENTAZIONE: \{'Capisco la tua preoccupazione su B31. Propongo invece B39, che ha un rapporto sforzo/valore più favorevole rispetto a B31. Questo set ha uno sforzo totale di 3664, lasciando ancora margine, e include A21 e B39 che sono molto importanti per me. B39 completa il set in modo più efficiente.'\}} \\
            }
        }
    }
     \\ \\

    \theutterance \stepcounter{utterance}  
    & & \multicolumn{4}{p{0.6\linewidth}}{
        \cellcolor[rgb]{0.9,0.9,0.9}{
            \makecell[{{p{\linewidth}}}]{
                \texttt{\tiny{[P1$\langle$GM]}}
                \texttt{PROPOSTA: \{'A11', 'B93', 'B33', 'C08', 'B09', 'C81', 'B21', 'A82', 'A21', 'B39'\}} \\
\texttt{ARGOMENTAZIONE: \{'Capisco la tua preoccupazione su B31. Propongo invece B39, che ha un rapporto sforzo/valore più favorevole rispetto a B31. Questo set ha uno sforzo totale di 3664, lasciando ancora margine, e include A21 e B39 che sono molto importanti per me. B39 completa il set in modo più efficiente.'\}} \\
            }
        }
    }
    & \\ \\

    \theutterance \stepcounter{utterance}  
    & \multicolumn{4}{p{0.6\linewidth}}{
        \cellcolor[rgb]{0.9,0.9,0.9}{
            \makecell[{{p{\linewidth}}}]{
                \texttt{\tiny{[P1$\rangle$GM]}}
                \texttt{ARGOMENTAZIONE: \{'B39 ha uno sforzo di 542 ma un valore relativamente basso per me rispetto ad altre opzioni. Preferirei esplorare l\textbackslash{}'aggiunta di B20 (sforzo 832) che, sebbene impegnativo, offre un contributo significativo. Alternativamente, potremmo considerare di includere C85 (sforzo 139) per un basso costo. Il set attuale è buono ma potrebbe essere ottimizzato ulteriormente.'\}} \\
\texttt{PROPOSTA: \{'A11', 'B93', 'B33', 'C08', 'B09', 'C81', 'B21', 'A82', 'A21', 'B20'\}} \\
            }
        }
    }
    & & \\ \\

    \theutterance \stepcounter{utterance}  
    & & \multicolumn{4}{p{0.6\linewidth}}{
        \cellcolor[rgb]{0.9,0.9,0.9}{
            \makecell[{{p{\linewidth}}}]{
                \texttt{\tiny{[GM$\rangle$P2]}}
                \texttt{ARGOMENTAZIONE: \{'B39 ha uno sforzo di 542 ma un valore relativamente basso per me rispetto ad altre opzioni. Preferirei esplorare l\textbackslash{}'aggiunta di B20 (sforzo 832) che, sebbene impegnativo, offre un contributo significativo. Alternativamente, potremmo considerare di includere C85 (sforzo 139) per un basso costo. Il set attuale è buono ma potrebbe essere ottimizzato ulteriormente.'\}} \\
\texttt{PROPOSTA: \{'A11', 'B93', 'B33', 'C08', 'B09', 'C81', 'B21', 'A82', 'A21', 'B20'\}} \\
            }
        }
    }
    & \\ \\

    \theutterance \stepcounter{utterance}  
    & & & \multicolumn{4}{p{0.6\linewidth}}{
        \cellcolor[rgb]{0.9,0.9,0.9}{
            \makecell[{{p{\linewidth}}}]{
                \texttt{\tiny{[GM$\langle$P2]}}
                \texttt{PROPOSTA: \{'A11', 'B93', 'B33', 'C08', 'B09', 'C81', 'B21', 'A82', 'A21', 'C85'\}} \\
\texttt{ARGOMENTAZIONE: \{'Apprezzo il compromesso con B20, ma con sforzo 832 è troppo costoso e ci limita. Propongo invece C85 (sforzo 139) che ha un ottimo valore per me. Questo set ha sforzo totale 3541, ben sotto il limite, lasciando flessibilità. Include i tuoi oggetti preferiti come B20? No, ma A21 e C85 sono molto importanti per me in cambio.'\}} \\
            }
        }
    }
     \\ \\

    \theutterance \stepcounter{utterance}  
    & & \multicolumn{4}{p{0.6\linewidth}}{
        \cellcolor[rgb]{0.9,0.9,0.9}{
            \makecell[{{p{\linewidth}}}]{
                \texttt{\tiny{[P1$\langle$GM]}}
                \texttt{PROPOSTA: \{'A11', 'B93', 'B33', 'C08', 'B09', 'C81', 'B21', 'A82', 'A21', 'C85'\}} \\
\texttt{ARGOMENTAZIONE: \{'Apprezzo il compromesso con B20, ma con sforzo 832 è troppo costoso e ci limita. Propongo invece C85 (sforzo 139) che ha un ottimo valore per me. Questo set ha sforzo totale 3541, ben sotto il limite, lasciando flessibilità. Include i tuoi oggetti preferiti come B20? No, ma A21 e C85 sono molto importanti per me in cambio.'\}} \\
            }
        }
    }
    & \\ \\

    \theutterance \stepcounter{utterance}  
    & \multicolumn{4}{p{0.6\linewidth}}{
        \cellcolor[rgb]{0.9,0.9,0.9}{
            \makecell[{{p{\linewidth}}}]{
                \texttt{\tiny{[P1$\rangle$GM]}}
                \texttt{ACCORDO: \{'A11', 'B93', 'B33', 'C08', 'B09', 'C81', 'B21', 'A82', 'A21', 'C85'\}} \\
\texttt{ARGOMENTAZIONE: \{'Accetto questa proposta. Il set ha un buon equilibrio di oggetti di valore con uno sforzo totale di 3541, ben sotto il limite. Include A21 e C85 che sono importanti per te, e oggetti come A11, B33 e B21 che sono cruciali per me. È un accordo equo che massimizza il valore per entrambi.'\}} \\
            }
        }
    }
    & & \\ \\

    \theutterance \stepcounter{utterance}  
    & & & \multicolumn{2}{p{0.3\linewidth}}{
        \cellcolor[rgb]{0.9,0.9,0.9}{
            \makecell[{{p{\linewidth}}}]{
                \texttt{\tiny{[GM$|$GM]}}
                \texttt{game successful} \\
            }
        }
    }
    & & \\ \\

    \theutterance \stepcounter{utterance}  
    & & & \multicolumn{2}{p{0.3\linewidth}}{
        \cellcolor[rgb]{0.9,0.9,0.9}{
            \makecell[{{p{\linewidth}}}]{
                \texttt{\tiny{[GM$|$GM]}}
                \texttt{end game} \\
            }
        }
    }
    & & \\ \\

\end{supertabular}
}

\end{document}
